\begin{ledgroupsized}[r]{120mm}
\footnotesize
\pstart
\noindent\textbf{Überlieferung:}
\pend
\end{ledgroupsized}
\begin{ledgroupsized}[r]{114mm}
\footnotesize
\pstart
\parindent -6mm
\makebox[6mm][l]{\textit{L}}%
Aufzeichnung:
LH~XXXVII~3 Bl.~63.
Ein Blatt~8\textsuperscript{o};
\href{\Scan{LH_37_03_063wz}}%
{Wasserzeichenfragment} in der Blattecke;
oberer und rechter Rand beschnitten.
Einzweidrittel Seiten.
\pend
\end{ledgroupsized}
\selectlanguage{latin}
\vspace{8mm}
\pstart
\centering
\href{\Scan{LH_37_03_063r}}%
{\lbrack63~r\textsuperscript{o}\rbrack} 
\hfill
Regula aequilibrii\indexS{regula aequilibrii}
\hfill
\phantom{\lbrack63~r\textsuperscript{o}\rbrack} 
\pend
\begin{edfigure}%
\includegraphics[scale=1]{%
gesamttex/edit_VIII-4/images/LH_37_03_063_d1_063r.pdf}
\edcaption[0.2em]{\lbrack\textit{Fig.~1}\rbrack}
\end{edfigure}
\pstart 
\noindent
\textit{CF} aequ 
\edtext{2\textit{CE}. \textit{AEB}}{%
\lemma{2\textit{CE}.}%
\Bfootnote{%
\textit{(1)}~\textit{AB} aequ.
\textit{(2)}~\textit{AEB}~\textit{L}}}
aequ. 2\textit{GFD}.
Primum \textit{ECF} axis aequilibrii\indexS{axis aequilibrii} est, nam ab utraque parte tantundem ponderis. 
\edtext{Ducatur \textit{LCN}}{\lemma{Ducatur}\Bfootnote{%
\textit{(1)}~semirecto angulo %
\textit{(2)}~\textit{LCN}~\textit{L}}}
recta angulo \textit{LCE}
\edtext{semirecto quae}{%
\lemma{semirecto}%
\Bfootnote{%
\textit{(1)}~itemque \textit{EL} et \textit{LB}
\textit{(2)}~quae~\textit{L}}}
sit novus axis librationis.\indexS{axis librationis} Ajo et hanc esse axem aequilibrii.\indexS{axis aequilibrii} 
Nam \textit{EL} et \textit{LB} 
\edtext{aequilibrantur quia per}{\lemma{aequilibrantur}\Bfootnote{%
\textit{(1)}~.~Sed et \textit{AHE} et \textit{GFD} quia rectae %
\textit{(2)}~.~Per
\textit{(3)}~quia per~\textit{L}}} 
rectae \textit{EB} centrum gravitatis\indexS{centrum gravitatis} 
\edtext{\textit{L}}{\lemma{}\Bfootnote{\textit{L} \textit{erg.~L}}} 
transit axis \textit{LN}. Similiter etiam \textit{AE} et \textit{GD} aequilibrantur nam \textit{FN} aequ. \textit{CH}. distantiae scilicet centrorum gravitatis\indexS{distantia centri gravitatis ab axe aequilibrii} \textit{H}. \textit{F} ab axe aequilibrii.\indexS{axis aequilibrii} 
Idem applicetur et ad reliquas proportiones. Idem scilicet si non de linea agatur sed de corporibus 
\edtext{ex centro grav.\indexS{centrum gravitatis}}{\lemma{}\Bfootnote{ex centro grav. \textit{erg.~L}}} 
suspensis in \textit{A}. \textit{H}. \textit{E}. \textit{L}. \textit{B}. 
et locis quibuslibet intermediis modo aequaliter utrobique.
\pend 
\begin{edfigure}%
\includegraphics[scale=1]{%
gesamttex/edit_VIII-4/images/LH_37_03_063_d2_063r.pdf}
\edcaption{\lbrack\textit{Fig.~2}\rbrack}
\end{edfigure}
\pstart 
\edtext{Demonstrabo semel quod duarum 
\edtext{linearum quae}{\lemma{linearum}\Bfootnote{%
\textit{(1)}~centr. gr. %
\textit{(2)}~quae~\textit{L}}}
rationem duplam habent centrum gravitatis\indexS{centrum gravitatis} 
sit in rectae earum centra gravitatis\indexS{centrum gravitatis} jungentis 
\edtext{intersectione simili,}{\lemma{intersectione}\Bfootnote{%
\textit{(1)}~in tres partes %
\textit{(2)}~simili, \textit{L}}}
videtur idem de caeteris etiam proportionibus hoc modo demonstrari.}{%
\lemma{Demonstrabo \bracks{...} demonstrari}%
\Cfootnote{Im wenig später entstandenen und thematisch eng zusammenhängenden Konzept N.\;\ref{57036} beweist Leibniz ein "Theorema aequiponderantium", das für alle Proportionen gilt und nicht nur für den Spezialfall von Linien, "quae rationem duplam habent".}}
\pend
\pstart
\textit{GD} et \textit{AL} centrum grav. in \textit{C}. Ponatur rursus $PA+LQ$ duplum $GD+AL$ enimque centr. g. $GD+AL$ sit \textit{C}. et centr. gr. $PA+LQ$ sit \textit{E}. 
Erit \textit{EC} rursus similiter secanda. Et ita res semper procedit repetitis sectionibus cum \textit{GD} est ad
\edtext{\textit{PQ} ut 1.}{\lemma{}\Bfootnote{%
\textit{PQ} \bars{est \textit{streicht Hrsg.}}
\textit{(1)}~in progressio %
\textit{(2)}~ut 1.~\textit{L}}}
ad
\href{\Scan{LH_37_03_063v}}%
{\lbrack63~v\textsuperscript{o}\rbrack}
numerum progressionis geometricae duplae. Sin non ita est, tunc pars ipsius \textit{GD}
\edtext{ut \textit{SD}}{\lemma{}\Bfootnote{ut \textit{SD} \textit{erg.~L}}}
sumenda est quae talis sit; et residui \textit{GS} rursum pars quae talis sit ad totum $SD+PQ$, ut 1. ad numerum progressionis geometricae, et ita 
\edtext{porro, et ita}{\lemma{porro,}\Bfootnote{%
\textit{(1)}~itaque accedemus %
\textit{(2)}~et ita~\textit{L}}}
semper procedendo demonstratum est in omni proportione quod
\edtext{de dupla}{\lemma{de}\Bfootnote{%
\textit{(1)}~geometrica %
\textit{(2)}~demonstra %
\textit{(3)}~dupla~\textit{L}}}
ostendimus.
\pend
\clearpage
\vspace{0.5em} %
\pstart
\noindent
\lbrack%
\textit{Nachfolgend kleingedruckter Text samt Zeichnung gestrichen:}\rbrack %
\pend
\nobreak
\vspace{0.5em}
\footnotesize
\pstart 
\noindent
Alia methodus 
\edlabel{37_03_063_1a}%
\edtext{}{%
\xxref{37_03_063_1a}{37_03_063_1b}%
\lemma{}\Bfootnote{demonstrandi \bars{brevissima \textit{gestr.}}. Sint~\textit{L}}}%
demonstrandi.
\pend
\normalsize
\begin{edfigure}%
\includegraphics[scale=1]{%
gesamttex/edit_VIII-4/images/LH_37_03_063_d3_063v.pdf}
\edcaption{\lbrack\textit{Fig.~3, gestrichen}\rbrack%
\pagelabel{Fig:37_03_063:Fig_3}}
\end{edfigure}
\footnotesize
\pstart 
Sint\edlabel{37_03_063_1b} duae rectae quaecunque \textit{AD} 
\edtext{et \textit{CB} axi aequilibrii\indexS{axis aequilibrii} \textit{EF} quaesito parallelae. Jung\lbrack at\rbrack ur}{\lemma{et \textit{CB}}\Bfootnote{%
\textit{(1)}~aequilibrii axi %
\textit{(a)}~parallelis
\textit{(b)}~parallelae,
\textit{(2)}~axi %
\textit{(a)}~aequilibrii \textit{EF}
\textit{(b)}~aequilibrii \textit{EF} \textit{EF} quaesito parallelae. %
\textit{(aa)}~Manifestum est in iisdem distantiis pondera %
\textit{(bb)}~\bars{Jungantur \textit{ändert Hrsg.}} \textit{GH}~\textit{L}}}
qui erit notus aequilibrii axis,\indexS{axis aequilibrii} secans quaesitum 
\edtext{in \textit{L}. Tamdiu}{\lemma{in \textit{L}.}\Bfootnote{%
\textit{(1)}~Ubi patet semper esse \textit{GL} ad \textit{LH}, ut (\textit{G})(\textit{L}) ad (\textit{L})(\textit{H}). et ita porro, adeoque semper eandem esse 
\textit{(a)}~distantiam ab axe
\textit{(b)}~distantiarum a centro grav. rationem %
\textit{(2)}~Tamdiu~\textit{L}}}
moveatur \textit{CB} sursum et deorsum, donec contingat \textit{GH}, \textit{EF} et \textit{DC} se secare in eodem puncto (\textit{L}). Hoc enim fieri potest.
\pend
\pstart
\textit{PH} ad \textit{KG} ut (\textit{H})\textit{M} ad (\textit{G})\textit{N}
{\normalsize \lbrack\textit{Text bricht ab.}\rbrack} 
\pend 
\normalsize
